\begin{itemize}\setlength\itemsep{0.1em}
\item \textbf{More data and larger $L$.} Add the complete $L=13$ set (94\,MB compressed) and the
$14\le L\le18$ samples. Because those samples are non-uniform, we will evaluate them per loop order
and report how error grows with $L$, which is the regime where cheap predictors matter most.
\item \textbf{Extrapolation.} Make training on low $L$ and testing on high $L$ a primary benchmark
(train $\le 12$, test 13; train $\le 13$, test 14--18), and study which features stay calibrated
across sizes (cycle counts grow with $L$ and may need per-$L$ normalization).
\item \textbf{Richer and more principled structural features.} Use longer cycle lengths and per-edge
cycle counts as edge features, sign- and basis-invariant Laplacian encodings (e.g.\ SignNet), and
subgraph-count features motivated by graph-theoretic invariants of periods (vertex cuts, Martin
invariants). We will also compare a pure ``set of node features'' model (DeepSets, no message
passing) to isolate what message passing adds.
\item \textbf{Ablations and tuning.} Sweep depth, width and pooling, train longer, use 5 seeds,
and examine per-length cycle-feature importance (dropping one cycle length at a time) as an
explainability analysis of which cycles determine the period.
\item \textbf{Cost.} Measure end-to-end wall-clock time per graph for our features + GNN against
the Hepp bound and against numerical integration, so that the accuracy/speed trade-off is explicit.
\item \textbf{Deliverables.} Port the pipeline (already modular and documented in \texttt{code/})
into a single Colab notebook and write the Medium blog post with these figures and PyG code snippets.
\end{itemize}

\paragraph{Code.} All code for this milestone is in \texttt{code/}: \texttt{data.py} (download
format, parsing, features, grouped splits), \texttt{models.py}, \texttt{train.py},
\texttt{baselines.py}, \texttt{make\_figures.py}, and shell scripts that reproduce every number in
this report.
